\chapter{Literature Survey}
\label{ch:litreview}

\section{Introduction}
A comprehensive review of existing literature reveals significant
research activity in AI-powered story generation, spanning traditional
rule-based approaches to modern deep-learning methods. This chapter
examines key contributions, identifies methodological trends and
highlights gaps that motivate the proposed research.

\section{Survey of Existing System}

\subsection{Evolution of Automated Story Generation}
Automated story generation has evolved through several distinct
phases. Early systems in the 1970s and 1980s relied on rule-based
approaches and story grammars that encoded narrative structures
explicitly. \emph{TALE-SPIN} \cite{meehan1977talespin} and
\emph{MINSTREL} demonstrated that computers could synthesise coherent
fables in a closed domain. The author-level goal-driven \emph{UNIVERSE}
system \cite{lebowitz1985universe} and the engagement-reflection
\emph{MEXICA} model \cite{perezperez2001mexica} extended these ideas
but remained constrained by handcrafted rules and narrow domain
coverage.

The statistical era brought data-driven approaches that learned
patterns from corpora of existing stories. While these systems showed
improved flexibility, they struggled with long-range coherence and
producing genuinely creative content. The neural network revolution
introduced recurrent architectures (LSTM, GRU) that captured
sequential dependencies more effectively but still suffered from
context degradation.

The transformer architecture introduced in 2017
\cite{vaswani2017attention} fundamentally changed the landscape of
natural-language generation. Attention mechanisms captured long-range
dependencies more effectively than recurrent approaches, leading to
dramatic improvements in generated-text quality. The subsequent
development of large-scale pretrained models such as
GPT-2 \cite{radford2019gpt2}, GPT-3 \cite{brown2020gpt3} and
GPT-4 \cite{openai2023gpt4} demonstrated that language models trained
on massive datasets could generate remarkably coherent and creative
narratives.

\subsection{Modern Interactive Story Generation Systems}
Several commercial and academic systems have explored interactive
storytelling: \emph{AI Dungeon} \cite{walton2019aidungeon} popularised
freeform LLM-driven interactive fiction; \emph{Dramatron}
\cite{mirowski2023dramatron} structures generation hierarchically
(log-line $\rightarrow$ scenes $\rightarrow$ dialogue) to improve
coherence; \emph{TaleBrush} \cite{chung2022talebrush} proposes
sketch-based co-creation. Recent CHI and IUI work
\cite{liu2023endusers,yuan2022wordcraft} examines mixed-initiative
authoring, finding that scaffolded interfaces improve novice usability
even when raw LLM quality is comparable.

\subsection{Multimodal Co-creation}
Text-to-image diffusion
\cite{ramesh2022unclip,rombach2022ldm,podell2024sdxl,bfl2024flux}
now produces high-quality scene art conditioned on natural language.
Recent work integrates image generation into narrative loops
\cite{schuhmann2024visualstory,bouhuis2024crossmodal}, demonstrating
reader-engagement uplift when illustrations are aligned with current
narrative state. On the audio side, browser-native speech synthesis
\cite{w3c2024webspeech} and neural TTS \cite{casanova2024xtts} permit
low-latency narration without paid APIs. Genre-conditioned ambient
music has been explored in game research \cite{lopezrincon2024music}
but is less studied for procedural narrative.

\subsection{Summary Table of Key Research}
Table~\ref{tab:literature} summarises ten representative works in
the field and identifies the gap each leaves unaddressed with respect
to interactive Q\&A-driven personalisation.

\begin{table}[!t]
\caption{Summary of recent literature in AI-driven story generation.}
\label{tab:literature}
\centering
\renewcommand{\arraystretch}{1.15}
\small
\begin{tabular}{@{}c c l p{4.8cm} p{2.2cm}@{}}
\toprule
\textbf{\#} & \textbf{Yr} & \textbf{Authors} & \textbf{Method} & \textbf{Gap} \\
\midrule
1 & 2024 & Farrell et al. \cite{farrell2024narrative}      & GPT heuristic search planner & No Q\&A \\
2 & 2023 & Huang et al. \cite{huang2023knowledge}          & GPT + knowledge graph        & Limited personalisation \\
3 & 2023 & Tambwekar et al. \cite{tambwekar2023controllable}& RL-controlled LSTM/LLM       & Not interactive \\
4 & 2023 & Li et al. \cite{li2023plotguided}               & Prompt-engineered GPT-3      & No Q\&A component \\
5 & 2022 & Ammanabrolu et al. \cite{ammanabrolu2022hier}   & Hierarchical Transformer     & No Q\&A branching \\
6 & 2022 & Mao et al. \cite{mao2022character}              & Character-centric Transformer & No Q\&A personalisation \\
7 & 2024 & Zhang et al. \cite{zhang2024prompt}             & Instruction-tuned LLM        & Single-shot prompts \\
8 & 2023 & Peng et al. \cite{peng2023dialogue}             & BART dialogue converter      & No structured Q\&A \\
9 & 2022 & Li et al. \cite{li2022discourse}                & Discourse planner + LLM      & Not personalised \\
10& 2024 & Wu et al. \cite{wu2024evaluation}               & LLM-based story metrics      & Evaluation only \\
\bottomrule
\end{tabular}
\end{table}

\subsection{Detailed Analysis of Selected Works}
\textbf{Farrell et al. (2024)} \cite{farrell2024narrative} propose
using large language models as heuristic guides for narrative
planning search algorithms. The system combines traditional AI
planning techniques with LLM-generated suggestions to create coherent
narrative sequences. While the approach produces well-structured
stories, it operates entirely without user interaction during the
generation process, limiting personalisation to initial parameter
settings.

\textbf{Huang et al. (2023)} \cite{huang2023knowledge} integrate
external knowledge graphs with GPT-based generation to improve
factual consistency and world-building in generated narratives. The
knowledge integration enables stories with richer detail and fewer
factual contradictions; however, the system treats users as passive
recipients with no mechanism for capturing user preferences through
interaction.

\textbf{Tambwekar et al. (2023)} \cite{tambwekar2023controllable}
introduce reinforcement learning techniques to control stylistic and
emotional aspects of generated stories. Users can specify desired
attributes like sentiment, writing style and thematic elements;
despite these controls, the generation process remains non-interactive
after initial specification, preventing dynamic adjustment based on
user feedback during story development.

\section{Limitations of Existing System / Research Gap}
\label{sec:gap}

\subsection{Gaps in User Interaction Models}
Farrell et al. (2024) propose LLM-based narrative planning but do
not support real-time Q\&A interaction, limiting dynamic story
adaptation based on user input. Huang et al. (2023) integrate
external knowledge graphs but treat users as passive participants
with no conversational involvement; they lack iterative feedback
mechanisms in which user answers influence future story context
and narrative flow.

\subsection{Gaps in Control and Personalisation}
Tambwekar et al. (2023) enable control over style and emotion but do
not allow users to shape the plot through interactive questioning.
Li et al. (2023) use static prompt-based plot guidance that cannot
adapt dynamically through user questions or answers, and do not
support dialogue-based refinement.

\subsection{Gaps in Narrative Structure}
Ammanabrolu et al. (2022) employ hierarchical neural planning but
lack user-centric interaction during story generation, following a
linear narrative structure without Q\&A-driven branching or decision
points. Existing approaches rely on predefined templates or
constraints for personalisation rather than live user responses;
template-based personalisation cannot capture the nuanced preferences
revealed through conversation.

\subsection{Gaps in Multimodal Augmentation}
None of the surveyed systems offer integrated illustration, narration
and ambient-music modalities that are dynamically conditioned on the
current segment text and detected genre. Existing image-augmented
storytelling efforts \cite{schuhmann2024visualstory,bouhuis2024crossmodal}
are research prototypes rather than deployable applications and do
not address the rate-limit problem of free image APIs.

\subsection{Gaps in End-User Authoring}
No existing storytelling system permits the user to add
\emph{entirely new questions} alongside the system's elicitation
protocol. Custom user-added Q\&A pairs that flow into the LLM prompt
through a dedicated context channel constitute a novel personalisation
pattern that this dissertation introduces.

\subsection{Summary of Gap Analysis}
The identified gaps converge on a central theme: existing systems
fail to position users as \emph{active co-creators} of narratives.
While technical capabilities for high-quality text generation exist,
the frameworks for leveraging these capabilities through meaningful,
multimodal, end-user-extensible interaction remain underdeveloped.
This dissertation addresses this fundamental gap with a comprehensive
Q\&A-based framework that enables genuine human-AI collaboration in
storytelling.

\section{Summary}
The literature survey reveals that although the components of
LLM-based generation, multimodal synthesis and conversational
interaction exist independently, no system combines all four into a
single accessible, provider-agnostic platform. The next chapter
formalises the problem statement, lists the project's objectives and
defines the work's scope.
